\section{Experiments}
\label{sec:experiments}

\subsection{Setup}
\textbf{Benchmark.} OmniDocBench v1.6~\cite{omnidocbench}: 1651 page images spanning books, textbooks, exam papers, magazines, academic literature, notes, newspapers, research reports and slides, in English and Chinese. We run the \emph{official} evaluation harness unmodified (quick-match matcher; text and reading order scored by normalized edit distance, display formulas by CDM~\cite{cdm} compiled with TeX~Live, tables by TEDS~\cite{teds}). The composite follows the leaderboard convention $\mathrm{Overall}=\big((1-\mathrm{TextEdit})\cdot 100+\mathrm{CDM}+\mathrm{TEDS}\big)/3$. We additionally report the v1.5 subset cut for comparability with rows published against v1.5.

\textbf{Hardware.} All PRISM measurements: a 16-logical-core consumer Windows 11 machine, CPU only, onnxruntime; no quantization beyond TATR INT8. Latency is end-to-end per page including all workers; peak RAM is the maximum working set of the whole process tree.

\subsection{Main results}
\label{sec:main_results}

\begin{table*}[t]
\centering
\caption{OmniDocBench v1.6\_full (1651 pages), official harness. Leaderboard rows from the official repository; PRISM measured by us on the identical dataset, harness and metrics. Abbreviated: the full 32-row table is in the supplementary.}
\label{tab:v16}
\small
\begin{tabular}{llccccc}
\toprule
Model & Type & Overall$\uparrow$ & Text$\downarrow$ & CDM$\uparrow$ & TEDS$\uparrow$ & Order$\downarrow$ \\
\midrule
MinerU2.5-Pro & VLM (1.2B) & 95.75 & 0.036 & 97.45 & 93.42 & 0.120 \\
PaddleOCR-VL & VLM (0.9B) & 94.18 & 0.040 & 95.91 & 90.65 & 0.135 \\
Gemini 3 Pro & General VLM & 92.91 & 0.064 & 95.99 & 89.15 & 0.165 \\
dots.ocr & VLM (3B) & 90.77 & 0.048 & 89.95 & 87.18 & 0.138 \\
GPT-5.2 & General VLM & 86.59 & 0.114 & 88.21 & 82.95 & 0.193 \\
MinerU-Pipeline & Pipeline (GPU) & 86.47 & 0.055 & 83.07 & 81.88 & 0.153 \\
\textbf{PRISM (ours)} & \textbf{Pipeline, CPU, 283\,MB} & \textbf{85.77} & 0.083 & 85.39 & 80.21 & 0.162 \\
olmOCR & VLM (7B) & 85.74 & 0.139 & 88.10 & 83.00 & 0.216 \\
Mistral OCR & VLM (API) & 85.66 & 0.097 & 89.91 & 76.78 & 0.171 \\
Nanonets-OCR-s & VLM (3B) & 83.61 & 0.108 & 81.46 & 80.18 & 0.213 \\
POINTS-Reader & VLM (3B) & 83.37 & 0.096 & 85.72 & 73.98 & 0.198 \\
Marker & Pipeline (GPU) & 78.44 & 0.157 & 85.24 & 65.77 & 0.243 \\
\bottomrule
\end{tabular}
\end{table*}

\Cref{tab:v16} places PRISM on the v1.6 leaderboard. PRISM lands above the 7B-parameter VLM class (olmOCR-7B, Mistral OCR), exceeds Marker --- a pipeline whose OCR alone is a 650M-parameter GPU VLM --- by 7.3 points, and sits within 0.7 of the MinerU GPU pipeline while beating its formula CDM outright (85.4 vs.\ 83.1), with every system above it using at least one order of magnitude more compute. By language, PRISM's text edit is 0.068 (EN) / 0.096 (ZH); table TEDS is 82.8 on Chinese pages and 75.6 on English ones (the residual English deficit is dense newspaper stock listings); the residual formula gap is Chinese-embedded math, bounded by the OCR-hybrid escape hatch (\cref{sec:formula}). On the v1.5 cut, PRISM scores \textbf{87.1} (text 0.073, CDM 85.3, TEDS 83.3) --- the \emph{highest-scoring pipeline} on that snapshot, ahead of PP-StructureV3 (86.7, which requires 8\,GB and 58\,s/page on CPU, \cref{tab:cpu_frontier}), with the same text edit distance and +1.6 TEDS; the systems above it are all GPU VLMs.

\subsection{Efficiency frontier}
\label{sec:efficiency}

\begin{table}[t]
\centering
\caption{Controlled CPU head-to-head: same 20 pages, same machine, isolated runs, 8-thread budget. PRISM row predates its current accuracy (conservative). $^{*}$format mismatch; $^{\dagger}$config artifact.}
\label{tab:cpu_frontier}
\small
\setlength{\tabcolsep}{3pt}
\begin{tabular}{lccccc}
\toprule
System & RAM & s/pg & Text$_{EN}\downarrow$ & Formula$_{EN}\downarrow$ & TEDS$\uparrow$ \\
\midrule
\textbf{PRISM} & \textbf{1.3\,GB} & \textbf{6.8} & 0.165 & 0.491 & 46.1 \\
PP-Struct.V3 (server) & 8.0\,GB & 58.2 & 0.081 & 0.320 & 56.5 \\
PP-Struct.V3 (mobile) & 3.6\,GB & 6.3 & 0.156 & 1.000$^{\dagger}$ & 43.5 \\
SmolDocling-256M & 2.3\,GB & 92.2 & 0.461 & 1.000$^{*}$ & 0.0$^{*}$ \\
GraniteDocling-258M & 2.3\,GB & 54.6 & 0.341 & 0.585 & 37.4 \\
\bottomrule
\end{tabular}
\end{table}

On the full-benchmark run PRISM's median latency is 5.6\,s/page (mean 6.6, p90 11.0) with a 2.60\,GB peak working set in the dual-worker benchmark configuration ($\sim$1.3\,GB single-worker); the speed-optimal v14 configuration (83.55 Overall) runs at 3.0\,s median, so the last +2.2 points cost +2.6\,s --- spent on a doubled formula decode budget and the higher-capacity layout model. \Cref{tab:cpu_frontier} is the same-machine comparison (run before the current build; PRISM's row there understates it): PRISM is simultaneously more accurate \emph{and} an order of magnitude faster than both sub-300M VLMs; only PP-StructureV3's server profile scored higher in that snapshot --- a gap the current build closes --- at $6\times$ the memory and $19\times$ the latency. Model inventory: layout 124\,MB, OCR 31\,MB, formulas 79\,MB, tables 7.4\,MB + 30\,MB TATR fallback $\approx$ 283\,MB.

\subsection{Recognition-verified normalization}
\label{sec:verified_exp}

\begin{table}[t]
\centering
\caption{Synthetic-defect study: mean per-block OCR edit distance (lower is better) on 40 text-heavy benchmark pages (20 EN / 20 ZH) under harsh photometric degradations with valid ground truth. \emph{none} = no correction, \emph{open} = open-loop camera stack, \emph{verified} = probe-gated stack.}
\label{tab:verified}
\small
\begin{tabular}{lccc}
\toprule
Input & none & open-loop & verified \\
\midrule
clean & \textbf{0.177} & 0.184 & 0.185 \\
shadow ($0.18\times$ + noise) & \textbf{0.181} & 0.185 & 0.188 \\
glare (saturating band) & \textbf{0.187} & 0.197 & 0.199 \\
tungsten cast & 0.187 & \textbf{0.183} & 0.185 \\
shadow + cast & \textbf{0.180} & 0.185 & 0.187 \\
\bottomrule
\end{tabular}
\end{table}

We evaluate the closed loop on photometric degradations of benchmark pages --- corner shadow to $0.18\times$ luminance with sensor noise, a saturating specular band, tungsten color cast, and their combination --- which keep ground-truth text-block geometry valid (\cref{tab:verified}). The study produced a finding we consider as valuable as the mechanism itself: \textbf{modern mobile OCR is largely invariant to photometric page defects}. A shadow that drops luminance to $0.18\times$ raises PP-OCRv6's block edit distance by only $0.004$; under every degradation, the \emph{uncorrected} image is at or near the best condition, and the open-loop correction stack pays a small tax everywhere (clean: $0.177 \rightarrow 0.184$). A page-level replication (whole-page reading against concatenated ground truth) agrees: verified gating tracks the no-correction baseline (0.600 vs.\ open-loop 0.603 on clean pages, 0.604 vs.\ 0.608 under glare) because the probe rejects proposals that do not help.

The implication inverts the conventional design: photometric normalization should not be a default enhancement stage but a \emph{verification-gated rescue} reserved for captures that recognition demonstrably cannot read as-is. Those cases are real but extreme --- on the 44 uncontrolled defect captures (hand shadows, tungsten lamps, page curl, deep corner shadows), the probe accepted shadow flattening on the 8 pages that pixel-statistics routing had misclassified as clean and \emph{rejected} it precisely where the divide erased strokes in dense paragraphs (a $7\times$ probe-score drop on one hand-shadow page); end-to-end, previously empty outputs (receipt-class captures) became fully readable. Open-loop stacks buy those rescues by taxing every clean page; the closed loop gets them for free.

\subsection{Ablations and negative results}
\label{sec:ablations}

\begin{table}[t]
\centering
\caption{Cumulative progression on the official v1.6 harness (identical dataset and scoring throughout). Each stage adds to the previous.}
\label{tab:ablation}
\small
\setlength{\tabcolsep}{4pt}
\begin{tabular}{lccccc}
\toprule
Configuration & Overall & Text$\downarrow$ & CDM & TEDS & Order$\downarrow$ \\
\midrule
Baseline (v9) & 70.37 & 0.158 & 56.90 & 69.96 & 0.323 \\
+ formula geometry, spans, rescues & 78.46 & 0.142 & 78.11 & 71.46 & 0.286 \\
+ answer-key rule, XY-cut, marginalia & 80.43 & 0.136 & 83.46 & 71.40 & 0.241 \\
+ SLANet-plus, OCRv6, sanitizer & 83.55 & 0.120 & 83.84 & 78.83 & 0.238 \\
+ DocLayoutV3 + native order, fml budget/hybrid & \textbf{85.77} & \textbf{0.083} & \textbf{85.39} & \textbf{80.21} & \textbf{0.162} \\
\bottomrule
\end{tabular}
\end{table}

\paragraph{What did not work.} We report failed directions with measurements, as they carry as much information as the successes:
\begin{itemize}
\item \textbf{Detector ensembling / dedicated formula detector.} Adding a YOLO-based math-formula detector raised CDM by 3.3 on formula pages but its false positives cost more text-edit than the CDM gained; net negative on the composite.
\item \textbf{Global confidence relaxation.} Lowering all detection gates to 0.15 raised recall but the marginal fragments stole matches from correct predictions; net negative. (Class-aware gates, \cref{sec:layout}, capture the benefit without the cost.)
\item \textbf{Frequency-domain moir\'e removal} erases text-line frequencies on paper photographs (\cref{sec:normalization}); enabled only for screen photographs.
\item \textbf{Emphasis markup.} Emitting bold/italic markers broke the matcher's caption pairing and turned matched blocks into unmatched 1.0-edit predictions; plain text scores strictly better.
\end{itemize}

\subsection{How much of the score is emission convention?}
\label{sec:perturb}

\begin{table}[t]
\centering
\caption{Emission-convention perturbation study: semantically-null transformations of our final predictions, re-scored with the unmodified official harness (1651 pages). $\Delta$ in points ($\times 100$).}
\label{tab:perturb}
\small
\setlength{\tabcolsep}{4pt}
\begin{tabular}{lcc}
\toprule
Perturbation (content unchanged) & $\Delta$Text & $\Delta$Order \\
\midrule
bold-wrap first sentence of every 3rd para & 0.00 & 0.00 \\
re-add GT marginalia verbatim & $+0.06$ & $-0.01$ \\
split paragraphs at sentence boundaries & $+0.41$ & $+1.54$ \\
\bottomrule
\end{tabular}
\end{table}

Because our development process kept surfacing evaluation-interaction effects, we quantified them systematically: we applied semantically-null transformations to our final predictions --- a human reader would extract identical content --- and re-scored with the untouched official harness (\cref{tab:perturb}). The harness proves \emph{robust} where content stays faithful: markdown styling is normalized away exactly, and even re-adding running headers/footers verbatim costs only 0.06 text points, because the matcher pairs them with the annotated-but-unscored ground-truth elements. The sensitivity is elsewhere: \textbf{segmentation granularity}. Emitting the same text as sentence-level blocks instead of paragraphs costs 0.41 text points and 1.54 reading-order points --- half a leaderboard step --- purely by convention. Combined with the interaction bugs we repaired during development (a line-break rewrite corrupting matrix row separators in 28\% of matched formula pairs; \emph{noisy OCR'd} marginalia, unlike verbatim ones, costing 0.36 text points through failed pairing; a single un-compilable macro zeroing a formula's CDM), the picture is consistent: benchmark scores conflate parsing ability with emission and segmentation conventions at the half-point scale, which is the same scale as many published leaderboard gaps. We recommend perturbation audits of this kind as standard practice when comparing document parsers.
